
Once the Knowledge Graph $\mathcal{G}$ is constructed, we transform the benchmark generation task into a path traversal problem. This ensures that every generated question is grounded in a verifiable logical chain. The structure of this component is shown in the Fig. \ref{fig:level3}.
% 一旦构建了知识图谱 $\mathcal{G}$，我们就将基准测试生成任务转化为路径遍历问题。这确保了生成的每个问题都基于可验证的逻辑链。

\textbf{Path-Based Query Synthesis.}
% 基于路径的查询合成。
We employ a Depth-First Search (DFS) algorithm to traverse $\mathcal{G}$, identifying directed paths $p = (e_{start}, r_1, e_1, \dots, e_{end})$. Each path represents a distinct causal or temporal sequence of medical events.
% 我们采用深度优先搜索（DFS）算法遍历 $\mathcal{G}$，识别有向路径 $p = (e_{start}, r_1, e_1, \dots, e_{end})$。每条路径代表一个不同的因果或时间序列的医疗事件。
For each extracted path, we retrieve the visual context and metadata of the nodes and edges. These structured signals are input into a generative MLLM (Qwen3-VL\cite{qwen3vl}) with a specialized prompt strategy:
% 对于提取出的每条路径，我们检索节点和边的视觉上下文和元数据。这些结构化信号被输入到具有专门提示策略的生成式多层逻辑模型（Qwen3-VL）中：
\begin{enumerate}
    \item \textbf{Input}: The structured path information (e.g., \textit{cotton swabs} $\xrightarrow{infiltration}$ \textit{iodine} $\xrightarrow{smear}$ \textit{skin}).
    \item \textbf{Instruction}: Synthesize a natural language question where $e_{start}$ serves as the premise and the temporal segment of $e_{end}$ serves as the answer.
    \item \textbf{Constraint}: The reasoning process must strictly follow the path topology.
\end{enumerate}
% 输入：结构化的路径信息（例如，棉签 浸润 碘伏 涂抹 皮肤）。
% 指令：生成一个自然语言问题，其中 $e_{start}$ 作为前提，时间片段 $e_{end}$ 作为答案。
% 约束：推理过程必须严格遵循路径拓扑结构。

\textbf{Adversarial Validation Mechanism.}
% 对抗性验证机制。
To ensure the rigorousness of the medical logic, we implement a critic-based validation loop. A separate, more powerful MLLM (GPT-4V\cite{GPT4V}) acts as an independent "Judge." 
% 为了确保医学逻辑的严谨性，我们实现了一个基于评判者的验证循环。一个独立的、功能更强大的机器学习模型（GPT-4V）作为独立的“评判者”。
It receives generated questions and original videos, and attempts to score and rank the questions by quality; only qualified questions are included in the dataset.
% 它接收生成的问题和原始视频，并尝试为这些问题进行打分与质量排序,只有合格的问题才会被纳入到数据集中.
This adversarial filtering minimizes the inclusion of ambiguous or unanswerable queries.
% 这种对抗性过滤机制最大限度地减少了包含歧义或无法回答的查询。